\scanpage{17}
\deleted{Supposons \unsure{quelque soit|quelles soient|qu'on soit} les corps résiduels de}

Notons qu'il y a dans $C$ un $S$-automorphisme canonique de symétrie, car $C$ est un schéma abélien sur $S$. Cet automorphisme \inserted{est d'ordre 2, il} opère sur les $\mathcal E_i$ [de façon compatible avec les structures multiplicatives \inserted{et les opérations de $\mathcal L$ sur $\mathcal E_\infty$}] et induit sur $\mathcal L$ le changt de signe. \deleted{Il en résulte aussitôt qu'il opère \unsure{trivialemt} dans $\mathcal E_2$. Et il y aura un \unsure{supplémentaire vers $\mathcal E_1$},} \deleted{Si les caractéristiques résiduelles de $S$ sont $\ne2$, on en conclut} Supposons les car. résiduelles de $S$ $\ne2$. Alors

\deleted{1) Il existe un supplémentaire unique de $\mathcal E_1$ dans $\mathcal E_2$ \unsure{stable} \unsure{qui sur lequel} $s$ opère}

% slanted
\margin{N.B. la \unsure{considération|construction|condition} de $C\to P(\mathcal E_2)$ et de la différent. de ce morphisme nous donne (sans restriction en car) un morphisme $\mathcal L^{\otimes3}\overset{\unsure{\sim}}{\longrightarrow}\mathcal E_3$ compatible avec les automorphismes (car fonctoriel en $S$).}

$s$ opère trivialemt dans $\mathcal E_2$ [c'est même sans la restriction qu'on vient de faire, on le verra], et il existe un supplémentaire unique \inserted{$L_3$} de $\mathcal E_2$ dans $\mathcal E_3$ où $s$ opère par symétrie, \deleted{\unsure{lequel}} qui est aussi égal à : $\mathcal L\cdot\mathcal E_2$, on le note $L_3$.

\deleted{Alors $L_3^{\otimes2}$ est un supplémentaire de $\mathcal E_5$ dans $\mathcal E_6$, soit $L_6$. A condition} On peut donc normaliser encore plus les bases admissibles de $\mathcal E_3$, en exigeant que $x$ soit une base de $\mathcal E_3$ ; donc le groupe des changemts de base admis est ceux de la forme
\[
\begin{pmatrix}1&\alpha&0\\0&\lambda^2&0\\0&0&\lambda^3\end{pmatrix}
\quad\text{i.e.}\quad
\begin{pmatrix}1&\unsure{\alpha|a}&0\\0&u&0\\0&0&v\end{pmatrix}
\quad\text{avec}\quad\deleted{\unsure{\lambda\lambda|\lambda}}u^3-v^2=0.
\]
Ce n'est autre que le \uline{groupe affine à 1 variable}.

Quelle est la normalisation correspondante de l'équation de $C$ ?
% archive sheet number: 16
